\documentclass[12pt]{article}

% Idioma e codificação
\usepackage[brazil]{babel}
\usepackage[utf8]{inputenc}
\usepackage[T1]{fontenc}

% Matemática e listas
\usepackage{amsmath, amssymb}
\usepackage{enumitem}

% Layout e cores
\usepackage{geometry}
\usepackage{graphicx}
\usepackage{xcolor}

% Diagramas e Portas Lógicas
\usepackage{tikz}
\usetikzlibrary{shapes.gates.logic.US, positioning, arrows.meta, calc, fit, shadows.blur}

% Cores institucionais
\definecolor{maincolor}{RGB}{0, 79, 151}
\definecolor{meuazul}{RGB}{0, 79, 151}

% Cabeçalho e rodapé
\usepackage{fancyhdr}

% Blocos
\usepackage[most]{tcolorbox}
\newtcolorbox{block}[1]{%
  colback=blue!5,
  colframe=maincolor,
  coltitle=white,
  title=#1,
  fonttitle=\bfseries,
  boxrule=1.5pt,
  arc=3mm,
  left=6pt,
  right=6pt,
  top=6pt,
  bottom=6pt
}

\newtcolorbox{resposta}[1]{%
  colback=white,
  colframe=meuazul!30,
  arc=1mm,
  height=#1,
  width=\textwidth,
  boxrule=0.8pt
}

% Margens
\geometry{margin=2.5cm}

% Fancyhdr
\pagestyle{fancy}
\fancyhf{}

\renewcommand{\headrulewidth}{2pt}
\renewcommand{\footrulewidth}{2pt}

\renewcommand{\headrule}{\hbox to\headwidth{%
  \color{meuazul}\leaders\hrule height \headrulewidth\hfill}}

\renewcommand{\footrule}{\hbox to\headwidth{%
  \color{meuazul}\leaders\hrule height \footrulewidth\hfill}}

\setlength{\headheight}{52pt}
\addtolength{\topmargin}{-20pt}

\fancyhead[L]{%
  \includegraphics[height=1.2cm]{logo.png}
}

\fancyhead[R]{%
  \raggedleft
  \textcolor{meuazul}{%
    \small
    \textbf{Lista de Exercícios — Desenvolvimento de Software para Dispositivos Móveis}\\
    Me.~Sergio Souza Novak\\
    \textit{Engenharia da Qualidade e Confiabilidade}
  }
}

\fancyfoot[C]{\textcolor{meuazul}{\thepage}}

% =========================
\begin{document}
% =========================

\vspace*{0.2cm}

\begin{center}
  {\Large \textbf{Conceitos iniciais android}}\\[0.2cm]
  \textit{Instruções: Justifique suas respostas de forma clara e objetiva.
  Para questões discursivas, escreva na caixa de resposta abaixo de cada enunciado.}
\end{center}

\vspace{0.4cm}

% ── Questão 1 ───────────────────────────────────────────────────────────────
\noindent\textbf{Questão 1 — Templates de Projeto}

\vspace{0.2cm}
\noindent
Ao iniciar um novo projeto no Android Studio, o desenvolvedor deve escolher um template inicial. Segundo o material da aula, qual é a principal diferença entre selecionar o template \textbf{Empty Activity} e o \textbf{Empty Views Activity}?

\vspace{0.3cm}
\noindent\textit{Resposta:}
\begin{resposta}{3.5cm}
\end{resposta}

\vspace{0.6cm}
\noindent\rule{\textwidth}{0.8pt}
\vspace{0.3cm}

% ── Questão 2 ───────────────────────────────────────────────────────────────
\noindent\textbf{Questão 2 — Manifest e Ponto de Entrada}

\vspace{0.2cm}
\noindent
O arquivo \texttt{AndroidManifest.xml} descreve informações essenciais do aplicativo para o sistema Android. Explique qual é a função específica da tag \texttt{<intent-filter>} contendo as propriedades \texttt{MAIN} e \texttt{LAUNCHER} declarada dentro de uma \texttt{<activity>}.

\vspace{0.3cm}
\noindent\textit{Resposta:}
\begin{resposta}{3.5cm}
\end{resposta}

\vspace{0.6cm}
\noindent\rule{\textwidth}{0.8pt}
\vspace{0.3cm}

% ── Questão 3 ───────────────────────────────────────────────────────────────
\noindent\textbf{Questão 3 — Estrutura de Recursos}

\vspace{0.2cm}
\noindent
Um projeto Android organiza seus códigos e recursos em diretórios separados. Dentro da pasta \texttt{res/}, encontramos diversas subpastas importantes. Cite e descreva a finalidade das pastas \texttt{drawable}, \texttt{layout} e \texttt{values}.

\vspace{0.3cm}
\noindent\textit{Resposta:}
\begin{resposta}{4.5cm}
\end{resposta}

\vspace{0.6cm}
\noindent\rule{\textwidth}{0.8pt}
\vspace{0.3cm}

% ── Questão 4 ───────────────────────────────────────────────────────────────
\noindent\textbf{Questão 4 — Boas Práticas: Centralização de Textos e Cores}

\vspace{0.2cm}
\noindent
Em vez de escrever textos ou definir códigos de cores diretamente nos layouts XML e no código Kotlin, o recomendado é usar os arquivos \texttt{strings.xml} e \texttt{colors.xml}. Quais são as principais vantagens dessa abordagem (por exemplo, em relação a traduções)?

\vspace{0.3cm}
\noindent\textit{Resposta:}
\begin{resposta}{3.5cm}
\end{resposta}

\vspace{0.6cm}
\noindent\rule{\textwidth}{0.8pt}
\vspace{0.3cm}

% ── Questão 5 ───────────────────────────────────────────────────────────────
\noindent\textbf{Questão 5 — View e ViewGroup}

\vspace{0.2cm}
\noindent
A interface gráfica de um aplicativo construído com Views obedece a uma estrutura hierárquica baseada nos conceitos de \textbf{View} e \textbf{ViewGroup}. Explique a diferença entre esses dois conceitos e dê pelo menos dois exemplos de cada um.

\vspace{0.3cm}
\noindent\textit{Resposta:}
\begin{resposta}{4.5cm}
\end{resposta}

\vspace{0.6cm}
\noindent\rule{\textwidth}{0.8pt}
\vspace{0.3cm}

% ── Questão 6 ───────────────────────────────────────────────────────────────
\noindent\textbf{Questão 6 — Unidades de Medida (\texttt{dp} vs. \texttt{sp})}

\vspace{0.2cm}
\noindent
Ao definir as dimensões e espaçamentos no Android, utilizamos a unidade \texttt{dp} (Density-independent pixels), enquanto para textos a recomendação é usar \texttt{sp} (Scale-independent pixels). Por que não devemos usar \texttt{dp} para definir o tamanho de fontes da interface?

\vspace{0.3cm}
\noindent\textit{Resposta:}
\begin{resposta}{4cm}
\end{resposta}

\vspace{0.6cm}
\noindent\rule{\textwidth}{0.8pt}
\vspace{0.3cm}

% ── Questão 7 ───────────────────────────────────────────────────────────────
\noindent\textbf{Questão 7 — Configuração de Tamanho no Layout}

\vspace{0.2cm}
\noindent
Um desenvolvedor deseja que um \texttt{Button} ocupe 100\% da largura disponível pelo seu elemento pai (container), mas que a sua altura se ajuste estritamente ao tamanho do texto inserido nele, sem sobrar espaços em branco desnecessários. Quais valores devem ser atribuídos aos atributos \texttt{layout\_width} e \texttt{layout\_height}, respectivamente?

\vspace{0.3cm}
\noindent\textit{Resposta:}
\begin{resposta}{3.5cm}
\end{resposta}

\vspace{0.6cm}
\noindent\rule{\textwidth}{0.8pt}
\vspace{0.3cm}

% ── Questão 8 ───────────────────────────────────────────────────────────────
\noindent\textbf{Questão 8 — Complexidade da Hierarquia de Views}

\vspace{0.2cm}
\noindent
Durante o desenvolvimento das telas de um aplicativo, é fundamental atentar-se à profundidade da hierarquia de componentes. Por que é considerado uma má prática de design de layout criar estruturas excessivamente profundas (com dezenas de ViewGroups aninhados)?

\vspace{0.3cm}
\noindent\textit{Resposta:}
\begin{resposta}{4cm}
\end{resposta}

\vspace{0.6cm}
\noindent\rule{\textwidth}{0.8pt}
\vspace{0.3cm}

% ── Questão 9 ───────────────────────────────────────────────────────────────
\noindent\textbf{Questão 9 — Build Gradle}

\vspace{0.2cm}
\noindent
O arquivo \texttt{build.gradle.kts} (nível de módulo da aplicação) desempenha um papel central na organização e compilação do projeto Android. Descreva, pelo menos, três tipos de configurações ou declarações que podem ser encontradas nesse arquivo e a sua finalidade (por exemplo: \texttt{defaultConfig}, \texttt{dependencies}, etc.).

\vspace{0.3cm}
\noindent\textit{Resposta:}
\begin{resposta}{4.5cm}
\end{resposta}

\vspace{0.6cm}
\noindent\rule{\textwidth}{0.8pt}
\vspace{0.3cm}

% ── Questão 10 ───────────────────────────────────────────────────────────────
\noindent\textbf{Questão 10 — Emulador Android}

\vspace{0.2cm}
\noindent
Para executar e testar o aplicativo sem um dispositivo físico, podemos usar o Emulador Android. Defina o que é esse emulador e explique o que é a ``System Image'' (Imagem do Sistema) e por que ela precisa ser baixada antes de iniciar um dispositivo virtual no Device Manager.

\vspace{0.3cm}
\noindent\textit{Resposta:}
\begin{resposta}{4cm}
\end{resposta}

\vspace{0.6cm}
\noindent\rule{\textwidth}{0.8pt}
\vspace{0.5cm}

\begin{center}
  \textit{``Boa atividade! Att. Me. Prof. Sergio Souza Novak.''}
\end{center}


% =========================
\end{document}
% =========================
